\section{模型改进与推广}

\subsection{模型改进方向}

\begin{enumerate}[label=（\arabic*）]
\item 将GAMM中的正态分布假设替换为Beta-GLMM。Y染色体游离DNA浓度取值限定在$(0,1)$区间且呈右偏分布，Beta分布能自然地建模这一有界连续变量。当前模型残差峰度4.15提示正态假设在尾部拟合不足，Beta-GLMM有望改善尾部预测精度，尤其是低浓度区域（接近4\%阈值）的概率估计。本文曾尝试拟合Beta-GLMM，但在高BMI子组上遇到收敛困难，未能得到稳定解，后续可通过调整初始值或采用贝叶斯采样器解决这一问题。

\item 将问题4的异常检测从单次截面分析扩展为纵向轨迹分析。对同一胎儿的多次检测记录，追踪Z值随孕周的变化轨迹，利用轨迹偏离模式（如斜率异常、波动异常）作为判别特征。纵向信息可能提供单次截面分析无法捕捉的动态异常信号，有望提升当前AUC=0.574的检测水平。

\item 利用临床结局数据（实际漏诊和延迟诊断记录）校准风险函数中的权重参数$w_d$和$w_a$。当前权重基于假设设定，灵敏度分析已表明高BMI组对权重选择高度敏感。通过回顾性临床数据建立延迟成本与浓度不足成本的经验比值，可使权重校准有据可依。

\item 采用贝叶斯层次模型替代频率学派的GAMM框架。贝叶斯方法通过先验分布对小样本子组实现部分池化（partial pooling），在高BMI组（仅20人）和极端BMI组（仅4人）的参数估计中引入合理的正则化，缩小当前蒙特卡洛分析中观察到的极宽置信区间（G4组标准差4.75周）。
\end{enumerate}

\subsection{推广应用}

\begin{enumerate}[label=（\arabic*）]
\item 风险函数优化框架适用于其他生物标志物驱动的产前筛查项目。例如PAPP-A（妊娠相关血浆蛋白A）和$\beta$-hCG（人绒毛膜促性腺激素）的检测同样存在浓度随孕周变化、检测时机影响准确性的问题。将风险函数中的浓度预测模型替换为对应标志物的GAMM，即可生成针对特定筛查项目的最优时点方案。

\item 残差分解技术可推广至涉及人体测量学变量共线性的其他临床建模场景。体重、BMI、体表面积、腰围等指标之间普遍存在强相关，直接纳入模型会导致VIF膨胀和参数不可解释。将衍生指标对基础指标回归取残差的策略，在药代动力学、营养流行病学等领域同样适用。

\item 蒙特卡洛误差传播方法可应用于其他需要在模型输出端量化不确定性的医学决策问题。当优化目标依赖于多步模型预测（如从GAMM预测到风险函数到最优时点），解析方法难以追踪误差传递路径时，蒙特卡洛模拟提供了一种通用的数值方案。
\end{enumerate}
